\chapter{Estructura compleja interna y geometría proyectiva de los estados de \(2\) niveles}
\label{chap:geometria-proyectiva-bloch}
\index{raíz interna de menos uno}
\index{Paley!matriz de conferencia}
\index{Witt}
\index{esfera de Bloch}
\index{M\"obius!acci\'on}
\index{grupo icosaédrico binario}

La geometría de fase empleada por la realización cuántica posee una
genealogía algebraica anterior a la notación compleja. Su punto de partida es
la incidencia finita de Paley--Witt; de ella se obtiene un endomorfismo cuyo
cuadrado es menos la identidad. La realificación conserva ese operador como
generador de fase, y la proyectivización de un espacio de dos niveles conduce
a la esfera de Bloch. Sobre esa esfera, la acción espinorial de
\(\mathrm{SU}(2)\), su expresión por transformaciones de Möbius y el grupo
icosaédrico binario forman realizaciones sucesivas de una misma estructura.

La cadena causal desarrollada en este capítulo es
\begin{equation}
 \boxed{
 \begin{gathered}
 C_P\longrightarrow A_W^2=-I_6,\\
 A_W^2=-I_6\longrightarrow
 \mathbb F_3^6\simeq\mathbb F_9^3
 \longrightarrow (E,J_E),\\
 (E,J_E)\longrightarrow
 \mathbb P(\mathcal H_J)\simeq\mathbb{CP}^{1}
 \longrightarrow S^2,\\
 S^2\longrightarrow
 \bigl\{\mathrm{SU}(2),\ \mathop{\text{M\"ob}}(S^2),\ 2A_5\bigr\}.
 \end{gathered}}
 \label{eq:cadena-paley-bloch-anexo}
\end{equation}
Cada flecha declara su dominio. La estructura cuadrática finita, la
realificación ortogonal y la geometría proyectiva permanecen relacionadas con
sus cuerpos de escalares explícitamente separados.

\section{La incidencia de Paley como origen algebraico de la fase}

Sea \(C_P\in M_6(\mathbb Z)\) la matriz simétrica de conferencia
\begin{equation}
 C_P=
 \begin{pmatrix}
 0& 1& 1& 1& 1& 1\\
 1& 0& 1&-1&-1& 1\\
 1& 1& 0& 1&-1&-1\\
 1&-1& 1& 0& 1&-1\\
 1&-1&-1& 1& 0& 1\\
 1& 1&-1&-1& 1& 0
 \end{pmatrix}.
 \label{eq:paley-conferencia-anexo}
\end{equation}
Sus filas satisfacen
\begin{equation}
 C_P C_P^{\mathsf T}=5I_6.
 \label{eq:paley-ortogonalidad-anexo}
\end{equation}
Al reducir módulo tres obtenemos
\begin{equation}
 A_W:=C_P\bmod 3\in M_6(\mathbb F_3).
 \label{eq:paley-reduccion-anexo}
\end{equation}

\begin{theorem}[Raíz interna de menos la identidad]
\label{thm:raiz-interna-menos-uno-anexo}
El endomorfismo \(J_W:v\mapsto vA_W\) de
\(V_W=\mathbb F_3^6\) satisface
\begin{equation}
 J_W^2=-I_{V_W}.
 \label{eq:Jw-cuadrado-anexo}
\end{equation}
En consecuencia, \(V_W\) admite una estructura canónica de módulo sobre
\(\mathbb F_3[X]/(X^2+1)\) una vez fijada la matriz de Paley.
\end{theorem}

\begin{proof}
La matriz \(C_P\) es simétrica. Reduciendo
\eqref{eq:paley-ortogonalidad-anexo} módulo tres se obtiene
\[
 A_W^2=5I_6\bmod3=2I_6=-I_6.
\]
Por ello, la evaluación \(X\mapsto J_W\) anula el ideal
\((X^2+1)\) y define la acción del cociente indicado.
\end{proof}

El polinomio \(X^2+1\) es irreducible sobre \(\mathbb F_3\), pues sus
valores en \(0,1,2\) son \(1,2,2\). Por tanto,
\begin{equation}
 \mathbb F_9:=\mathbb F_3[\iota]/(\iota^2+1)
 \label{eq:F9-desde-paley-anexo}
\end{equation}
es un cuerpo y \(J_W\) realiza la multiplicación por \(\iota\).

\begin{corollary}[Descenso de seis coordenadas ternarias a tres coordenadas cuadráticas]
\label{cor:F3seis-F9tres-anexo}
La regla
\begin{equation}
 (a+b\iota)v:=av+bJ_Wv,
 \qquad a,b\in\mathbb F_3,
 \label{eq:accion-F9-anexo}
\end{equation}
convierte \(V_W\) en un espacio vectorial de dimensión tres sobre
\(\mathbb F_9\). En las coordenadas fijadas por \(A_W\),
\begin{equation}
 \mathbb F_3^6\simeq\mathbb F_9^3.
 \label{eq:F3seis-igual-F9tres-anexo}
\end{equation}
\end{corollary}

\begin{proof}
La identidad \(J_W^2=-I\) verifica las reglas de multiplicación de
\(\mathbb F_9\). Como
\([\mathbb F_9:\mathbb F_3]=2\) y
\(\dim_{\mathbb F_3}V_W=6\), la dimensión sobre \(\mathbb F_9\) es tres.
\end{proof}

El mismo operador organiza el código ternario mediante el grafo
\begin{equation}
 C_W=\{(q,qA_W):q\in\mathbb F_3^6\}\subset\mathbb F_3^{12}.
 \label{eq:codigo-grafo-J-anexo}
\end{equation}
Si \(G=(I_6\mid A_W)\), entonces
\[
 GG^{\mathsf T}=I_6+A_W A_W^{\mathsf T}=0.
\]
Así, la raíz interna de menos la identidad participa simultáneamente en la
estructura cuadrática y en la isotropía del código. Esta doble función explica
por qué la fase y la incidencia excepcional comparten un antecedente y
conservan lectores diferentes.

\section{Realificación ortogonal de la fase}

La fase analítica se construye sobre un espacio real antes de elegir una carta
compleja. Sea \((E,g)\) un espacio euclídeo real de dimensión par y sea
\(J_E\in\operatorname{End}(E)\) un operador tal que
\begin{equation}
 J_E^2=-I,
 \qquad
 g(J_Eu,J_Ev)=g(u,v).
 \label{eq:estructura-compleja-ortogonal-anexo}
\end{equation}
La segunda identidad implica \(J_E^{\mathsf T}=-J_E\).

\begin{proposition}[Fase como grupo ortogonal de un parámetro]
\label{prop:fase-ortogonal-J-anexo}
Para todo \(\theta\in\mathbb R\),
\begin{equation}
 R_J(\theta):=e^{\theta J_E}
 =\cos\theta\,I+\sin\theta\,J_E
 \label{eq:rotacion-J-anexo}
\end{equation}
es ortogonal y satisface
\begin{equation}
 R_J(\theta)R_J(\phi)=R_J(\theta+\phi),
 \qquad R_J(2\pi)=I.
 \label{eq:grupo-fase-J-anexo}
\end{equation}
\end{proposition}

\begin{proof}
La serie exponencial se separa en potencias pares e impares y, usando
\(J_E^2=-I\), produce la expresión trigonométrica. La antisimetría de
\(J_E\) da
\(R_J(\theta)^{\mathsf T}=R_J(-\theta)\); por tanto,
\(R_J(\theta)^{\mathsf T}R_J(\theta)=I\). La ley de composición resulta de
que todos los factores son funciones del mismo operador.
\end{proof}

La elección de la carta \(J_E\leftrightarrow i\) convierte \(E\) en un
espacio vectorial complejo: para \(a,b\in\mathbb R\),
\begin{equation}
 (a+ib)u:=au+bJ_Eu.
 \label{eq:carta-compleja-J-anexo}
\end{equation}
En esa carta, \(R_J(\theta)\) se escribe como multiplicación por
\(e^{i\theta}\). La notación compleja publica la fase que el operador real ya
transporta.

\begin{remark}[Dos escalas, una misma relación cuadrática]
El operador \(J_W\) actúa sobre \(\mathbb F_3^6\); el operador \(J_E\) actúa
sobre un espacio real con producto escalar. La prolongación
\(J_W\rightsquigarrow J_E\) es una realización que conserva la relación
\(J^2=-I\) y declara además topología, norma y completación. Este orden evita
usar la estructura analítica continua como premisa de la incidencia finita.
\end{remark}

\section{La recta proyectiva compleja y la esfera de Bloch}

Sea \(\mathcal H_J\) un espacio de Hilbert de dimensión compleja dos obtenido
mediante la carta \eqref{eq:carta-compleja-J-anexo}. Un estado puro normalizado
es un vector
\begin{equation}
 \psi=\begin{pmatrix}z_0\\ z_1\end{pmatrix},
 \qquad |z_0|^2+|z_1|^2=1,
 \label{eq:espinor-normalizado-anexo}
\end{equation}
y su estado físico proyectivo es el rayo
\([\psi]=\{e^{i\vartheta}\psi:\vartheta\in\mathbb R\}\).
El espacio de rayos es \(\mathbb P(\mathcal H_J)\simeq\mathbb{CP}^{1}\).

Definimos las coordenadas
\begin{equation}
 \begin{aligned}
 r_1&=2\operatorname{Re}(\overline z_0z_1),\\
 r_2&=2\operatorname{Im}(\overline z_0z_1),\\
 r_3&=|z_0|^2-|z_1|^2.
 \end{aligned}
 \label{eq:coordenadas-bloch-anexo}
\end{equation}

La representación esférica de los estados de dos niveles se formula en las
coordenadas introducidas por Bloch \cite{Bloch1946}.

\begin{theorem}[Identificación proyectiva de Bloch]
\label{thm:CP1-esfera-bloch-anexo}
La aplicación
\begin{equation}
 \mathcal B:\mathbb{CP}^{1}\longrightarrow S^2,
 \qquad [\psi]\longmapsto(r_1,r_2,r_3),
 \label{eq:mapa-bloch-anexo}
\end{equation}
es una biyección suave. El proyector de rango uno asociado a \([\psi]\) es
\begin{equation}
 \rho_\psi=|\psi\rangle\langle\psi|
 =\frac12\bigl(I_2+r_1\sigma_1+r_2\sigma_2+r_3\sigma_3\bigr),
 \label{eq:proyector-bloch-anexo}
\end{equation}
donde \(\sigma_1,\sigma_2,\sigma_3\) son las matrices de Pauli.
\end{theorem}

\begin{proof}
La multiplicación de \(\psi\) por una fase común deja invariantes las tres
coordenadas. Además,
\[
 r_1^2+r_2^2+r_3^2
 =4|z_0|^2|z_1|^2+(|z_0|^2-|z_1|^2)^2=1.
\]
Toda \(\boldsymbol r\in S^2\) determina el proyector positivo
\(\rho=(I+\boldsymbol r\cdot\boldsymbol\sigma)/2\), que satisface
\(\rho^2=\rho\) y \(\operatorname{tr}\rho=1\); su imagen es un único rayo.
La fórmula matricial directa de \(|\psi\rangle\langle\psi|\) proporciona
\eqref{eq:proyector-bloch-anexo}. Las cartas afines usuales de
\(\mathbb{CP}^{1}\) hacen suave el mapa y su inversa.
\end{proof}

La esfera de Bloch registra la clase proyectiva y conserva, a través de
\(\rho_\psi\), todas las esperanzas de observables hermíticos de dos niveles.
La fase global pertenece a la fibra \(S^1\) del fibrado de Hopf
\cite{Hopf1931}
\begin{equation}
 S^1\longrightarrow S^3\longrightarrow S^2.
 \label{eq:hopf-bloch-anexo}
\end{equation}
La distinción HMT entre estado enriquecido y coordenada publicada encuentra
aquí una realización precisa: el vector normalizado pertenece a \(S^3\), el
lector proyectivo publica un punto de \(S^2\), y la fibra conserva la fase que
la proyección identifica.

\section{Acción espinorial, rotaciones y transformaciones de Möbius}

Todo elemento de \(\mathrm{SU}(2)\) puede escribirse como
\begin{equation}
 U=
 \begin{pmatrix}
 \alpha&\beta\\
 -\overline\beta&\overline\alpha
 \end{pmatrix},
 \qquad |\alpha|^2+|\beta|^2=1.
 \label{eq:SU2-parametrizacion-anexo}
\end{equation}
Actúa linealmente sobre espinores y, por conjugación, sobre proyectores.

\begin{theorem}[Doble recubrimiento y acción en la esfera]
\label{thm:SU2-SO3-bloch-anexo}
Para cada \(U\in\mathrm{SU}(2)\) existe una única
\(R_U\in\mathrm{SO}(3)\) tal que
\begin{equation}
 U(\boldsymbol r\cdot\boldsymbol\sigma)U^\dagger
 =(R_U\boldsymbol r)\cdot\boldsymbol\sigma.
 \label{eq:accion-SU2-bloch-anexo}
\end{equation}
La aplicación \(U\mapsto R_U\) es un homomorfismo sobreyectivo con núcleo
\(\{\pm I_2\}\).
\end{theorem}

\begin{proof}
La conjugación conserva el espacio real tridimensional de matrices hermíticas
sin traza y su producto escalar
\(\langle A,B\rangle=\tfrac12\operatorname{tr}(AB)\). Por ello induce una
transformación ortogonal de \(\mathbb R^3\). La continuidad y el valor en la
identidad fijan determinante positivo. Si la conjugación es trivial sobre las
tres matrices de Pauli, \(U\) conmuta con toda matriz compleja \(2\times2\),
y por el lema de Schur es escalar; la condición de determinante uno deja
\(U=\pm I_2\). La representación eje--ángulo muestra que toda rotación posee
dos preimágenes opuestas.
\end{proof}

En la carta afín \(z=z_0/z_1\) de \(\mathbb{CP}^{1}\), la misma acción toma
la forma
\begin{equation}
 z\longmapsto
 \frac{\alpha z+\beta}
      {-\overline\beta z+\overline\alpha}.
 \label{eq:mobius-SU2-anexo}
\end{equation}
La rotación de Bloch y la transformación fraccional son, por tanto, dos
coordenadas de la misma acción proyectiva. La primera emplea el proyector
hermítico; la segunda, la esfera de Riemann.

\section{La especialización icosaédrica y el grupo binario \texorpdfstring{\(2A_5\)}{2A5}}

El grupo rotacional del icosaedro es un subgrupo de
\(\mathrm{SO}(3)\) isomorfo a \(A_5\). Su preimagen por el recubrimiento del
\cref{thm:SU2-SO3-bloch-anexo} define el grupo icosaédrico binario
\begin{equation}
 2A_5:=\pi^{-1}(A_5)\subset\mathrm{SU}(2).
 \label{eq:def-2A5-anexo}
\end{equation}

\begin{proposition}[Extensión central icosaédrica]
\label{prop:extension-central-2A5-anexo}
Existe una sucesión exacta
\begin{equation}
 1\longrightarrow\{\pm I_2\}
 \longrightarrow2A_5
 \longrightarrow A_5
 \longrightarrow1,
 \label{eq:sucesion-2A5-anexo}
\end{equation}
y \(|2A_5|=120\). Un giro icosaédrico de órdenes tres o cinco se levanta a
elementos de órdenes seis o diez, respectivamente.
\end{proposition}

\begin{proof}
La restricción de \(\pi:\mathrm{SU}(2)\to\mathrm{SO}(3)\) a la preimagen de
\(A_5\) es sobreyectiva y conserva el núcleo \(\{\pm I_2\}\). La exactitud y
el teorema del orden dan \(|2A_5|=2|A_5|=120\). Una rotación de ángulo
\(2\pi/n\) se representa en \(\mathrm{SU}(2)\) con semiángulo
\(\pi/n\); después de \(n\) iteraciones alcanza \(-I_2\) y después de
\(2n\), la identidad.
\end{proof}

La especialización icosaédrica fija una bisagra rigurosa entre la geometría de
estados de dos niveles, las transformaciones de Möbius y la rama de simetrías
finitas. El desarrollo de McKay y de las estructuras excepcionales utiliza
después la representación de \(2A_5\) y sus módulos; su contenido adicional
se desarrolla en los teoremas específicos de esa rama.

\section{Composición functorial \texorpdfstring{\(\mathbb F_3^6\to\mathbb F_9^3\to\mathbb{CP}^{1}\simeq S^2\to\mathrm{SU}(2)\to\mathrm{SO}(3)\to 2A_5\)}{F3^6 -> F9^3 -> CP1 -> S2 -> SU(2) -> SO(3) -> 2A5}}

La relación \(J^2=-I\) atraviesa cuatro categorías conservando su tipo:
\begin{enumerate}
 \item en \(\mathbb F_3^6\), define la extensión cuadrática
       \(\mathbb F_9\) y organiza la incidencia de Witt;
 \item en un espacio real ortogonal, genera el flujo de fase
       \(e^{\theta J_E}\);
 \item en dimensión compleja dos, la proyectivización produce
       \(\mathbb{CP}^{1}\simeq S^2\) y los proyectores de Bloch;
 \item la acción de \(\mathrm{SU}(2)\) desciende a rotaciones y a
       transformaciones de Möbius, y su preimagen icosaédrica produce
       \(2A_5\).
\end{enumerate}
La estructura compleja interna posee así una procedencia finita y una
realización analítica explícita. La esfera de Bloch publica la geometría de los
rayos; la genealogía HMT conserva además la ruta, la orientación, el carry y la
memoria que anteceden a esa publicación.
